% =====================================================================
\chapter{Referencial teórico}\label{cap:referencial}
% =====================================================================

O referencial teórico a seguir constrói a base conceitual que suporta a hipótese
de menor endividamento em empresas familiares, partindo das teorias clássicas e
aprofundando-se nas particularidades de sua governança.

% ---------------------------------------------------------------------
\section{Conflitos de agência e a riqueza socioemocional (SEW)}\label{sec:sew}
% ---------------------------------------------------------------------

O arcabouço teórico que define as empresas familiares enfatiza o envolvimento da
família no controle, propriedade e gestão \cite{santos2018}. Entre as companhias
analisadas nesta pesquisa, a Baumer ilustra a presença simultânea dessas três
dimensões: a família detém cerca de 89\% das ações ordinárias, integra o bloco de
controle e ocupa a presidência do conselho de administração e a
diretoria-presidência. Já a RaiaDrogasil exemplifica um arranjo comum no mercado
brasileiro, em que três famílias reunidas em um mesmo bloco de controle detêm, em
conjunto, menos da metade das ações ordinárias, mas mantêm membros no conselho de
administração.
\verificar{participação e cargos atuais na BALM e na RADL no Formulário de
Referência mais recente}

Em contraste com o modelo de agência clássico, que foca nos conflitos entre
gestores e acionistas, estudos sobre empresas familiares destacam que essas
organizações frequentemente perseguem objetivos que transcendem a maximização do
retorno econômico, incorporando aspectos relacionados à preservação do controle,
da identidade familiar e da continuidade do negócio ao longo das gerações
\cite{ames2020,santos2018}. Nesse contexto, ganha destaque o conceito de riqueza
socioemocional (SEW), que abrange elementos como a manutenção da identidade
familiar, a autonomia decisória e o desejo de perpetuidade do negócio
\cite{santos2018}.

\citeonline{santos2018}, ao analisarem a remuneração executiva, fornecem
evidências de que a prioridade em preservar o SEW pode se manifestar em incentivos
menos intensos em remuneração pecuniária, com remuneração total e variável
menores, sugerindo um alinhamento de interesses intrínseco que minimiza o conflito
clássico entre gestor e proprietário. Contudo, a prevalência do SEW gera um
conflito latente com os credores, uma vez que a dívida, por meio de mecanismos de
monitoramento e restrições contratuais, pode representar uma ameaça à manutenção
do controle e da autonomia da família controladora, aspectos considerados centrais
para a preservação da riqueza socioemocional \cite{santos2018,ames2020}.

% ---------------------------------------------------------------------
\section{Endividamento elevado e o risco de interferência}\label{sec:risco}
% ---------------------------------------------------------------------

A dívida corporativa, embora ofereça benefícios fiscais, eleva o risco de
dificuldade financeira. Um dos mecanismos contratuais mais relevantes nesse
contexto é o \textit{covenant}, uma cláusula restritiva presente em contratos de
dívida que estabelece limites para determinados indicadores
econômico-financeiros e atua como instrumento de monitoramento e proteção dos
credores \cite{costa2018}.

A relação entre endividamento e escolhas contábeis não é linear.
\citeonline{costa2018} encontraram evidências de que, ao atingir níveis muito
elevados de endividamento, as empresas brasileiras tendem a aumentar o componente
discricionário dos lucros para evitar a quebra de cláusulas contratuais. Nesse
contexto, a redução da qualidade dos lucros pode ser percebida pelos gestores como
menos onerosa do que os custos decorrentes da violação dos \textit{covenants},
incluindo a renegociação compulsória ou o vencimento antecipado das obrigações
financeiras. Essa relação de risco é especialmente associada às dívidas de longo
prazo, que normalmente apresentam cláusulas restritivas mais rigorosas.

O prazo da dívida, contudo, expõe a empresa aos credores de outra forma. Enquanto a
dívida de longo prazo concentra as restrições nas cláusulas contratuais, a dívida
de curto prazo precisa ser renovada com frequência e sujeita a empresa ao risco de
liquidez, isto é, ao risco de que os credores se recusem a refinanciá-la ou o
façam em condições desfavoráveis \cite{diamond1991}. Esse risco tende a se agravar
em períodos de crise, quando o crédito se torna mais escasso. Para uma família que
busca preservar a autonomia decisória, depender de renovações frequentes significa
submeter-se repetidamente à avaliação dos credores. Assim, além de um menor nível
de dívida, a aversão à interferência de credores pode se manifestar na preferência
por prazos mais longos.

Para as empresas familiares, o risco de violação de \textit{covenants} pode ser
particularmente relevante, uma vez que a eventual intervenção dos credores pode
limitar a autonomia decisória da família controladora. Considerando que a
preservação do controle e da independência da gestão constitui um dos principais
componentes da riqueza socioemocional, a exposição excessiva ao endividamento pode
representar uma ameaça a objetivos que transcendem o desempenho econômico da firma
\cite{santos2018,pestana2021}. Portanto, espera-se que a família controladora
apresente maior aversão à alavancagem excessiva como forma de reduzir o risco de
perda de controle.

Há, porém, um argumento em sentido oposto. Como a emissão de novas ações dilui a
participação da família, a preservação do controle pode levar a família a preferir
a dívida ao capital próprio. \citeonline{poletti2022}, em empresas
latino-americanas, encontraram empresas familiares mais alavancadas quando a
participação da família é menor, como forma de manter o controle, e redução da
alavancagem quando essa participação é elevada. No Brasil, \citeonline{kayo2018}
observaram empresas familiares mais endividadas e com ajuste mais lento da
estrutura de capital. Essas evidências ajudam a explicar por que a literatura
permanece inconclusiva. Ainda assim, considerando a perspectiva da riqueza
socioemocional e o papel dos credores discutido nesta seção, adota-se como
hipótese a direção de menor endividamento.
\verificar{conferir nos artigos de Poletti-Hughes e Martínez Garcia (2022) e de
Kayo, Brunaldi e Aldrighi (2018) os resultados descritos neste parágrafo}

Com base no exposto, formula-se a seguinte hipótese:

\begin{quote}
\textbf{H1: Empresas de controle familiar apresentam um nível de alavancagem
financeira inferior ao de empresas não familiares.}
\end{quote}
